%% cover_letter.tex — G6: Full VSM Dynamics (AVR + PSS Extension)
%% IEEE Transactions on Energy Conversion
\documentclass[11pt,a4paper]{letter}
\usepackage[T1]{fontenc}
\usepackage[utf8]{inputenc}
\usepackage[margin=2.5cm]{geometry}
\usepackage{microtype}

\signature{Anoop~V.~Eluvathingal\\
           Ph.D.\ Candidate, Dept.\ of Electrical Engineering\\
           Indian Institute of Technology Madras\\
           Chennai 600036, India\\[2pt]
           \texttt{arun.eluvathingal@iitm.ac.in}}

\address{Department of Electrical Engineering\\
         Indian Institute of Technology Madras\\
         Chennai 600036, India}

\begin{document}
\begin{letter}{Editor-in-Chief\\
               \textit{IEEE Transactions on Energy Conversion}}

\opening{Dear Editor,}

We submit for your consideration the manuscript:

\begin{center}
  \textbf{Extending the SAMBPS Estimation Window Beyond 150\,ms:
  A Nine-State VSM Model with Full AVR and PSS Dynamics}
\end{center}

\textbf{Background and motivation.}
Model-based adaptive protection of synchronous and virtual synchronous
machine (VSM) units requires accurate fault-current estimation over the
full protection operating interval.  Our prior work (under review, IEEE
Transactions on Power Systems) demonstrated that a sixth-order Park~$dq$
truth model limits the self-adaptive model-based protection system
(SAMBPS) to estimation windows of at most 150\,ms.  Beyond this
boundary, automatic voltage regulator (AVR) and power system stabiliser
(PSS) dynamics introduce unmodelled excitation transients that raise the
Levenberg--Marquardt condition number $\kappa_n$ above the SAMBPS veto
threshold, causing false-restrain outcomes for slow-clearing faults and
island-detection scenarios.

\textbf{Contributions.}
This paper makes four distinct contributions:

\begin{enumerate}
  \item \textbf{Nine-state VSM model (C1)}: A stiff ODE formulation
        coupling six Park~$dq$ flux states with one IEEE~Type-1 AVR
        excitation state and two lead-lag PSS states, with a
        virtual-impedance droop interface.

  \item \textbf{Observability analysis (C2)}: Proof that the
        nine-state model is locally observable from terminal
        $(V_t, I_t, \omega)$ measurements with $L^* \geq 40$\,ms,
        with explicit conditions on fault current and PSS gain.

  \item \textbf{Window extension proposition (C3)}: Formal proof
        that incorporating AVR and PSS states extends the
        $\kappa_n < 30$ operating region from $L^* \leq 150$\,ms
        to $L^* \leq 380$\,ms ($+230$\,ms, $+153\%$), subject to
        $\pm 10\%$ parameter knowledge.

  \item \textbf{Islanded-microgrid case study (C4)}: On a 10-bus
        benchmark, zone detection accuracy at 300\,ms improves
        from 84.0\% (6th-order) to 97.1\% (9-state), and all
        island-detection scenarios that the 6th-order model
        failed are correctly resolved.
\end{enumerate}

\textbf{Fit to the journal.}
\textit{IEEE Transactions on Energy Conversion} is the primary venue
for machine modelling, excitation control, and generator dynamics
studies.  This manuscript directly connects AVR and PSS dynamics to
protection system performance --- a bridge between the energy
conversion and power systems communities that fits squarely within
the journal's scope.

\textbf{Suggested reviewers.}
\begin{itemize}
  \item Prof.\ Q.-C.\ Zhong (Illinois Institute of Technology)
        — virtual synchronous machine control
  \item Prof.\ Marta Molinas (NTNU) — GFM inverter dynamics
  \item Prof.\ Dionysios Aliprantis (Purdue University)
        — machine modelling and simulation
  \item Dr.\ Silvio D'Arco (SINTEF Energy Research)
        — VSM stability and control
\end{itemize}

This manuscript is not under review elsewhere.  All authors have
approved the submission.

\closing{Sincerely,}
\end{letter}
\end{document}
